\documentclass[11pt]{article}

\usepackage[margin=1in]{geometry}
\usepackage{amsmath}
\usepackage{amssymb}
\usepackage{amsthm}
\usepackage{booktabs}
\usepackage[hidelinks]{hyperref}
\setlength{\emergencystretch}{3em}

\newtheorem{theorem}{Theorem}
\newtheorem{lemma}{Lemma}
\newtheorem{corollary}{Corollary}
\newtheorem*{proposition}{Proposition}
\newtheorem{propositionnum}{Proposition}

\theoremstyle{definition}
\newtheorem{definition}{Definition}

\theoremstyle{remark}
\newtheorem{remark}{Remark}
\newtheorem*{remarkunnum}{Remark}

\newcommand{\Z}{\mathbf{Z}}
\newcommand{\Q}{\mathbf{Q}}
\newcommand{\R}{\mathbf{R}}
\newcommand{\C}{\mathbf{C}}
\newcommand{\Fun}{\mathbf{F}_1}
\newcommand{\NS}{\mathrm{NS}}
\newcommand{\Hom}{\mathrm{Hom}}
\newcommand{\End}{\mathrm{End}}
\newcommand{\Corr}{\mathrm{Corr}}
\newcommand{\Pic}{\mathrm{Pic}}
\newcommand{\pr}{\mathrm{pr}}
\newcommand{\covol}{\mathrm{covol}}
\newcommand{\trdeg}{\mathrm{trdeg}}
\newcommand{\SpecZ}{\operatorname{Spec}\Z}

\title{Classical correspondences between per-prime Tate curves:\\
a matching obstruction}
\author{Kunal Tyagi\thanks{\texttt{hello.jay.tyagi@gmail.com}.
\emph{Use of AI.} The original work used Claude (Anthropic). This revision,
proof review, and companion checks used OpenAI Codex. The author directed the
work and is responsible for the content. Code and publication artifacts:
\url{https://github.com/x67ai/riemann-rh-program}.}}
\date{August 26, 2026; revised October 18, 2026 (version 2)}
\hypersetup{pdftitle={Classical correspondences between per-prime Tate curves: a matching obstruction},pdfauthor={Kunal Tyagi}}
\begin{document}
\maketitle
\begin{abstract}
For the complex Tate curves $E_p=\C^\times/p^\Z$ attached to primes, we compute
the classical divisorial correspondence groups. For distinct primes,
$\Hom(E_p,E_q)\ne0$ exactly when
$\log p\log q\in4\pi^2\Q$. If $\log p\log q/(4\pi^2)=u/v$ in lowest positive
terms, this group has rank one and its generator has isogeny degree $uv$.
The isogeny relation between distinct primes is a partial matching:
unconditionally, each prime has at most one partner. These rigidity statements
specialize results of Winkelmann; the explicit Hom generator and degree are
included here. Outside the matching, the N\'eron--Severi group of $E_p\times E_q$
consists of the two fiber classes, which act trivially on $H^1$.
On every diagonal surface $\End(E_p)=\Z$, by Gelfond--Schneider, and the
intersection profile of multiplication graphs is independent of $p$.
Consequently no connected classical correspondence system can involve three
distinct prime fibers. This is a restriction on cross-fiber assembly, not a
universal prohibition of prime-indexed endomorphism monoids or of
explicit-formula realizations with additional coefficient data. Version 2
withdraws that broader no-go assertion of version 1 and states the exact
support restriction required by the obstruction.
\end{abstract}

\section{Introduction and scope of the revision}\label{sec:summary}
Connes and Consani's construction supplies the complex elliptic curves
$E_p=\C^\times/p^\Z$ \cite{CC26}. Their ordinary products have classical
intersection theory. A natural question is how much of a correspondence
calculus can be assembled from these products. Here a correspondence from
$E$ to $E'$ means a divisor class on $E\times E'$, modulo the two fiber
classes. This specific classical category is the subject of the paper.

Theorems~\ref{thm:isogeny}--\ref{thm:diagonal} compute its morphisms and
intersection data. The rational-product criterion and matching argument
specialize Winkelmann's rigidity \cite{Wi04}, rather than constituting new
transcendence theorems. The off-diagonal rational-product question remains
open; the matching bound and diagonal non-CM result are unconditional.
Theorem~\ref{thm:nogo} gives the resulting connected-component obstruction.

\paragraph{Correction of version 1.}
Version 1, titled \emph{Products of the per-prime Tate curves of absolute
geometry carry no correspondence calculus for the Weil explicit formula},
claimed that its three proposed requirements excluded every assignment of
correspondence data. That claim is withdrawn. A multiplicative prime-indexed
monoid need not use maps between different fibers, and constancy of an
unmarked intersection lattice does not prohibit varying test coefficients.
The corrected theorem requires an explicitly connected cross-fiber system.
The consequences for more general explicit-formula proposals are discussed
in Section~\ref{sec:scope}; no universal impossibility result is asserted.

\section{Lattices and classical correspondences}\label{sec:background}
\subsection{The lattice presentation}
\label{sec:notation}

Throughout, $p, q, \dots$ denote primes; set
\[
y_p := \frac{\log p}{2\pi} > 0 , \qquad
\tau_p := i\, y_p , \qquad
\Lambda_p := \Z \oplus \Z\cdot\tau_p , \qquad
E_p := \C/\Lambda_p .
\]

\begin{lemma}[lattice presentation; {\cite[\S4.1.2]{CC26}}]
\label{lem:lattice}
$\exp_2\colon \C \to \C^\times$, $w \mapsto e^{2\pi i w}$, is a surjective
homomorphism with kernel $\Z$, and
\[
\exp_2(w + \tau_p) = e^{2\pi i w}\cdot e^{2\pi i \cdot (i y_p)}
 = e^{2\pi i w}\cdot e^{-\log p} = p^{-1}\cdot \exp_2(w) .
\]
Since $p^\Z = (p^{-1})^\Z$ as subgroups of $\C^\times$, $\exp_2$ induces an
isomorphism $\C/\Lambda_p \cong \C^\times/p^\Z$. \qed
\end{lemma}

The presentation was confirmed numerically to 55 digits (check~A of the
numerical suite; see \emph{Provenance}). The sign of $\tau_p$ is
immaterial: $\Z \oplus \Z(-\tau_p) = \Lambda_p$.

Standard facts used without further comment: every homomorphism of complex
tori $\C/\Lambda \to \C/\Lambda'$ is induced by a unique $\lambda \in \C$
with $\lambda\Lambda \subseteq \Lambda'$ (lift to the universal cover); every
holomorphic map between complex tori is a homomorphism followed by a
translation (rigidity); $\Hom(E, E')$ denotes the group of holomorphic
homomorphisms.

\section{The isogeny criterion}
\label{sec:isogeny}

Theorems~\ref{thm:isogeny}--\ref{thm:coherence} specialize to primes a
rigidity of Winkelmann~\cite[Lemma~7 and the proof of Thm.~2, p.~16]{Wi04},
who studies the curves $\C^\times/\lambda^{\Z}$ attached to eigenvalues
$\lambda$ of elements of a discrete subgroup of $\mathrm{SL}_2(\C)$. He shows
that isogeny of two members of that family forces
$4\pi^2/(\log\lambda_i \log\lambda_j) \in \Q$, divides two such relations and
appeals to multiplicative independence to conclude that each curve has at most
one isogenous partner, and runs the argument on three curves; his Lemma~7 is
stated as an equivalence, and the converse follows from it in one line, so
both directions of the criterion below are his. Proofs are given in full here
because the paper is self-contained and the prime-indexed forms differ in
presentation from his. What is added is the identification of the group
$\Hom(E_p, E_q)$, its rank, and the degree $uv$ of the minimal isogeny.

\begin{theorem}
\label{thm:isogeny}
Let $p \neq q$ be distinct primes. Then
\[
\Hom(E_p, E_q) \neq 0
\iff y_p\, y_q \in \Q
\iff \log p \cdot \log q \in 4\pi^2\Q .
\]
Moreover, when the condition holds and $y_p y_q = u/v$ in lowest terms
($u, v \in \Z_{>0}$), the full homomorphism group is
\[
\Hom(E_p, E_q) = \Z\cdot\lambda_0 , \qquad \lambda_0 = i\, v\, y_q ,
\]
of rank $1$, and the isogeny induced by $\lambda_0$ has degree $uv$.
\end{theorem}

\begin{proof}
A nonzero homomorphism is multiplication by some $\lambda \in \C^\times$ with
$\lambda\Lambda_p \subseteq \Lambda_q$. Two containment conditions:

\begin{enumerate}
\item $\lambda\cdot 1 \in \Lambda_q$: so
$\lambda = a + b\cdot\tau_q = a + i\, b\, y_q$ for some $a, b \in \Z$.
\item $\lambda\cdot\tau_p \in \Lambda_q$: compute
$\lambda\cdot\tau_p = (a + i b y_q)\cdot(i y_p) = -b\, y_p y_q + i\, a\, y_p$.
Writing this as $c + i\, d\, y_q$ with $c, d \in \Z$:
\begin{itemize}
\item real part: $-b\, y_p y_q = c \in \Z$;
\item imaginary part: $a\, y_p = d\, y_q$.
\end{itemize}
\end{enumerate}

\emph{Case $a \neq 0$.} Then $d \neq 0$ (as $y_p, y_q > 0$) and
$y_p/y_q = d/a$, i.e.\ $\log p/\log q \in \Q$, i.e.\
$a\cdot\log p = d\cdot\log q$ with $a, d$ nonzero integers, i.e.\
$p^a = q^d$ (after clearing signs) --- impossible for distinct primes by
unique factorization. So $a = 0$, hence $d = 0$.

\emph{Case $a = 0$.} Then $\lambda = i\, b\, y_q$ with $b \neq 0$ (else
$\lambda = 0$), and the surviving condition is $b\, y_p y_q \in \Z$, which
forces $y_p y_q \in \Q$. This proves the forward direction, and identifies
$\Hom = \{\, i\, b\, y_q : b \in \Z,\ b\, y_p y_q \in \Z \,\}$. With
$y_p y_q = u/v$ in lowest terms, $b\, u/v \in \Z \iff v \mid b$, so
$\Hom = \Z\cdot(i\, v\, y_q)$.

\emph{Converse.} If $y_p y_q = u/v$, then $\lambda_0 = i\, v\, y_q$ satisfies
$\lambda_0\cdot 1 = v\cdot\tau_q \in \Lambda_q$ and
$\lambda_0\cdot\tau_p = (i v y_q)(i y_p) = -v\, y_p y_q = -u \in \Z \subseteq
\Lambda_q$, so $\lambda_0\Lambda_p \subseteq \Lambda_q$ and
$\lambda_0 \neq 0$ --- a genuine isogeny. (Numerically confirmed on a
synthetic pair with $y\, y' = 3/7$ exactly --- check~B.)

\emph{Degree.} $\covol(\Lambda_p) = y_p$ and
$\covol(\lambda_0\Lambda_p) = |\lambda_0|^2\cdot y_p = v^2 y_q^2\cdot y_p$,
so $\deg = [\Lambda_q : \lambda_0\Lambda_p] = v^2 y_q^2 y_p / y_q
= v^2\cdot(y_p y_q) = v^2\cdot(u/v) = uv$.

Finally $y_p y_q \in \Q \iff \log p \log q = 4\pi^2\cdot y_p y_q \in 4\pi^2\Q$.
\end{proof}

\begin{remark}[bookkeeping reconciliation]
\label{rem:bookkeeping}
The criterion may be stated as $\log p\cdot\log q \in 4\pi^2\Q$ or derived in
the form $b\cdot\log p\cdot\log q = -4\pi^2 c$ and phrased as
$\log p\cdot\log q \in \pi^2\Q$. These are the \emph{same condition}:
$4\pi^2\Q = \pi^2\Q$ as subsets of $\R$, because $4 \in \Q^\times$ --- a
rational multiple of $\pi^2$ is a rational multiple of $4\pi^2$ and
conversely. The normalization-free statement is $y_p\cdot y_q \in \Q$; we
display the constant as $4\pi^2$ because
$y_p y_q = \log p\cdot\log q/(4\pi^2)$ makes it the natural one.
\end{remark}

\begin{remark}[symmetry]
\label{rem:symmetry}
The criterion is symmetric in $p$ and $q$; hence
$\Hom(E_p, E_q) \neq 0 \iff \Hom(E_q, E_p) \neq 0$ (as the general theory
also guarantees via the dual isogeny). ``Exceptional partner'' is therefore
a symmetric relation.
\end{remark}

\section{Two-partner rigidity and the three-prime consequence}
\label{sec:familykill}

\begin{theorem}[two-partner rigidity]
\label{thm:twopartner}
For every prime $p$ there is at most one prime $q \neq p$ with
$\Hom(E_p, E_q) \neq 0$. Consequently the set of ``exceptional pairs''
$\{p, q\}$ (pairs with nonzero $\Hom$), if nonempty, forms a \emph{partial
matching} on the primes: no prime appears in two of them.
\end{theorem}

\begin{proof}
Suppose $q \neq q'$ are two such primes. By Theorem~\ref{thm:isogeny},
$y_p y_q \in \Q^\times$ and $y_p y_{q'} \in \Q^\times$. Dividing:
$y_q / y_{q'} = (y_p y_q)/(y_p y_{q'}) \in \Q$, i.e.\
$\log q / \log q' \in \Q$, i.e.\ $q^m = q'^n$ for some positive integers
$m, n$ --- impossible for distinct primes by unique factorization.
\end{proof}

The proof is unique factorization plus linear algebra over $\Q$ applied to
the products $y_p y_q$, and it is agnostic about whether any exceptional
pair exists.

\begin{theorem}[the three-prime consequence]
\label{thm:coherence}
Among the three surfaces $E_2\times E_3$, $E_2\times E_5$, $E_3\times E_5$,
at most one carries a nonzero correspondence group; in particular no two of
them --- a fortiori not all three --- can simultaneously carry nonzero
correspondence groups. (By Theorem~\ref{thm:NS} below,
``correspondence group of $E_p\times E_q$'' $= \Hom(E_p, E_q)$.)
\end{theorem}

\begin{proof}
Any two of the three pairs $\{2,3\}, \{2,5\}, \{3,5\}$ share a prime, so two
simultaneously nonzero groups would give some prime two exceptional
partners, contradicting Theorem~\ref{thm:twopartner}. Directly: if
$y_2 y_3 = r_1 \in \Q^\times$ and $y_2 y_5 = r_2 \in \Q^\times$ then
$y_3/y_5 = r_1/r_2 \in \Q$, i.e.\ $\log 3/\log 5 \in \Q$, i.e.\
$3^m = 5^n$ --- impossible; the other two cases are identical.
\end{proof}

\section{N\'eron--Severi collapse off the matching}
\label{sec:NS}

We first fix the classical dictionary. For elliptic curves $E, E'$ over
$\C$, the group of \emph{divisorial correspondences} from $E$ to $E'$ is
$\NS(E \times E')$ modulo the subgroup generated by the two fiber classes
$\xi_1 = [\{0\}\times E']$, $\xi_2 = [E\times\{0\}]$ (Weil's ``trivial
correspondences''):
\[
\Corr(E, E') := \NS(E \times E') / \langle \xi_1, \xi_2 \rangle .
\]

\begin{theorem}
\label{thm:NS}
Let $E, E'$ be elliptic curves over $\C$. Then there is a canonical
isomorphism
\[
\NS(E \times E') \cong \Z\cdot\xi_1 \oplus \Z\cdot\xi_2 \oplus \Hom(E, E') ,
\]
hence $\Corr(E, E') \cong \Hom(E, E')$ (\cite[Ch.~5]{BL}; self-contained
proof below). Consequently, for distinct primes $p \neq q$ with
$\Hom(E_p, E_q) = 0$ (i.e., every pair outside the --- conjecturally
empty --- exceptional matching of
Theorems~\ref{thm:isogeny}--\ref{thm:twopartner}):
\begin{itemize}
\item[\textup{(a)}] \textbf{Rank collapse.}
$\NS(E_p \times E_q) = \Z\xi_1 \oplus \Z\xi_2$ has rank $2$, with
intersection form $\xi_1^2 = \xi_2^2 = 0$, $\xi_1\cdot\xi_2 = 1$: every
divisor class is numerically $a\xi_1 + b\xi_2$, and $\Corr(E_p, E_q) = 0$.
\item[\textup{(b)}] \textbf{$H^1$-inertness.} Every divisor class on
$E_p \times E_q$ acts as zero on $H^1$ (in both directions,
$H^1(E_p) \to H^1(E_q)$ and the transpose). The correspondence calculus is
not merely small; as a calculus of operators it is the zero calculus.
\item[\textup{(c)}] \textbf{No graphs, no diagonal.} The only holomorphic
maps $E_p \to E_q$ are constants, so the only graph classes are fiber
classes. There is no diagonal: not as a curve (a ``diagonal'' in a product
of nonisomorphic curves means the graph of an isomorphism, and none exists),
and not even as a numerical class --- no $D \in \NS$ satisfies the
diagonal's numerical triple $(D\cdot\xi_1, D\cdot\xi_2, D^2) = (1, 1, 0)$,
since $D\cdot\xi_1 = D\cdot\xi_2 = 1$ forces $D = \xi_1 + \xi_2$ and then
$D^2 = 2 \neq 0$.
\item[\textup{(d)}] \textbf{Castelnuovo--Severi vacuous.} The
Hodge-index/Castelnuovo--Severi inequality on a product of curves is the
statement that the intersection form is negative semidefinite on the
\emph{primitive} part $\langle \xi_1, \xi_2 \rangle^\perp \subset \NS$. Here
that primitive part is $0$: if $D = a\xi_1 + b\xi_2$ satisfies
$D\cdot\xi_1 = b = 0$ and $D\cdot\xi_2 = a = 0$ then $D = 0$. Equivalently:
every class of bidegree $(d_1, d_2)$ satisfies $D^2 = 2 d_1 d_2$ ---
Castelnuovo--Severi holds \emph{with equality everywhere}, carrying zero
information. The positivity the proposal was to import fiber-by-fiber is free
here because there is nothing left for it to constrain.
\end{itemize}
\end{theorem}

\begin{proof}[Proof of the decomposition]
All inputs are standard abelian-variety theory (\cite[Ch.~2]{BL} for the
Poincar\'e bundle, seesaw, and rigidity; Ch.~5 for the endomorphism context;
the statement itself is the classical correspondence dictionary). Write
$\hat{E}' = \Pic^0(E')$ for the dual, $P$ for the Poincar\'e bundle on
$\hat{E}' \times E'$, normalized so that
\[
P|_{\{\alpha\}\times E'} \cong \alpha\quad(\alpha \in \hat{E}'),
\qquad P|_{\hat{E}'\times\{0\}}\cong\mathcal O.
\]

\emph{Step 1 (the homomorphism $\varphi_L$).} For $L \in \Pic(E\times E')$
put $L_x := L|_{\{x\}\times E'}$; $\deg L_x$ is constant in $x$. Normalize:
\[
\widetilde{L} := L \otimes \pr_1^*\bigl(L|_{E\times\{0\}}\bigr)^{-1}
 \otimes \pr_2^*(L_0)^{-1} ,
\]
so that $\widetilde{L}_x = L_x \otimes L_0^{-1} \in \Pic^0(E')$ for all $x$,
$\widetilde{L}_0$ is trivial, and $\widetilde{L}|_{E\times\{0\}}$ is
trivial. By the universal property of $(\hat{E}', P)$ there is a unique
morphism $\varphi_L\colon E \to \hat{E}'$ with
$(\varphi_L \times \mathrm{id}_{E'})^* P \cong \widetilde{L}$ (the a-priori
ambiguity by $\pr_1$-pullbacks is killed by restricting to $E\times\{0\}$,
where both sides are trivial). Since $\widetilde{L}_0$ is trivial,
$\varphi_L(0) = 0$; a pointed morphism of abelian varieties is a
homomorphism (rigidity). So $\varphi_L \in \Hom(E, \hat{E}')$.

\emph{Step 2 (kernel; factorization through $\NS$).} $L \mapsto \varphi_L$
is additive: normalization and restriction are additive in $L$, and
uniqueness plus the biadditivity of the Poincar\'e bundle
($((\mathrm{add})\times\mathrm{id})^* P \cong \pr_{13}^* P \otimes
\pr_{23}^* P$ on $\hat{E}'\times\hat{E}'\times E'$, \cite[Ch.~2]{BL}) give
$\varphi_{L\otimes L''} = \varphi_L + \varphi_{L''}$. If
$L \in \Pic^0(E\times E')$ then, since the dual of a product is the product
of the duals, $L = \pr_1^* M \otimes \pr_2^* M'$ with $M \in \Pic^0(E)$,
$M' \in \Pic^0(E')$; the normalization annihilates it, so $\varphi_L = 0$
and the map factors through $\NS(E\times E')$. Conversely if
$\varphi_L = 0$ then $\widetilde{L}_x$ is trivial for every $x$ and
$\widetilde{L}|_{E\times\{0\}}$ is trivial, so by the seesaw principle
$\widetilde{L}$ is trivial; hence
\[
L \cong \pr_1^*\bigl(L|_{E\times\{0\}}\bigr) \otimes \pr_2^*(L_0) ,
\]
whose N\'eron--Severi class is
$\deg\bigl(L|_{E\times\{0\}}\bigr)\cdot\xi_1 + \deg(L_0)\cdot\xi_2$ (the
class of $\pr_1^*$ of a degree-$d$ bundle is
$d\cdot[\{\mathrm{pt}\}\times E'] = d\cdot\xi_1$, and symmetrically). So the
kernel of $\NS \to \Hom(E, \hat{E}')$ is exactly $\Z\xi_1 \oplus \Z\xi_2$ (a
free rank-$2$ group: the intersection numbers against $\xi_2, \xi_1$ read
off the coefficients).

\emph{Step 3 (surjectivity and splitting).} For $\psi \in \Hom(E, \hat{E}')$
set $L^\psi := (\psi \times \mathrm{id}_{E'})^* P$. Then
$L^\psi_x = P|_{\{\psi(x)\}\times E'} = \psi(x) \in \Pic^0(E')$, $L^\psi_0$
is trivial, and $L^\psi|_{E\times\{0\}} = \psi^*\bigl(P|_{\hat{E}'\times\{0\}}\bigr)$
is trivial --- so $L^\psi$ is already normalized and
$\varphi_{L^\psi} = \psi$ by uniqueness. Additivity of
$\psi \mapsto [L^\psi]$ is the Poincar\'e biadditivity again. This splits
the sequence:
\[
\NS(E\times E') \cong \Z\xi_1 \oplus \Z\xi_2 \oplus \Hom(E, \hat{E}') .
\]

\emph{Step 4 (eliminating the dual).} For an elliptic curve,
$x \mapsto [\mathcal{O}_{E'}([x] - [0])]$ is a canonical isomorphism
$E' \cong \hat{E}'$ (the Abel--Jacobi isomorphism),
so $\Hom(E, \hat{E}') \cong \Hom(E, E')$.
\end{proof}

\begin{proof}[Proof of \textup{(a)}--\textup{(d)}]
(a) is immediate from the decomposition with $\Hom = 0$, plus
$\xi_1^2 = \xi_2^2 = 0$ (disjoint parallel fibers) and
$\xi_1\cdot\xi_2 = 1$ (one transverse point).

(b) The action of a correspondence on cohomology,
$\alpha \mapsto \pr_{2*}([D] \cup \pr_1^*\alpha)$, depends only on the class
$[D] \in H^2(E\times E', \Z)$. K\"unneth:
$H^2 = (H^2\otimes H^0) \oplus (H^1\otimes H^1) \oplus (H^0\otimes H^2)$. A
class $\pr_1^*\beta$ ($\beta \in H^2(E)$) acts on $\alpha \in H^1(E)$ by
$\pr_{2*}(\pr_1^*(\beta \cup \alpha)) = 0$ since
$\beta \cup \alpha \in H^3(E) = 0$; a class $\pr_2^*\beta$ acts by
$\beta \cup \pr_{2*}(\pr_1^*\alpha) = 0$ since fiber integration drops the
degree of $\pr_1^*\alpha \in H^1$ below zero. Now $\xi_1 = \pr_1^*[\mathrm{pt}]$
and $\xi_2 = \pr_2^*[\mathrm{pt}]$, so with $\NS = \Z\xi_1 \oplus \Z\xi_2$
every \emph{algebraic} class in $H^2$ has vanishing
$H^1\otimes H^1$-component and acts as zero on $H^1$. The transpose
direction is identical using $\Hom(E_q, E_p) = 0$
(Remark~\ref{rem:symmetry}).

(c) Rigidity: every holomorphic map $E_p \to E_q$ is a homomorphism followed
by a translation; $\Hom = 0$ leaves the constants, whose graphs are
horizontal fibers $E_p \times \{c\} \equiv \xi_2$. The numerical computation
is displayed in the statement.

(d) The calculation in the statement gives the orthogonal complement
of the fiber classes directly. No use of a Riemann hypothesis is involved.
The general Hodge index theorem is not needed to prove that this particular
primitive subgroup is zero.
\end{proof}

\begin{remark}[what still exists on these surfaces]
\label{rem:stillexists}
Rank-$2$ $\NS$ does not mean the surfaces are curve-poor: $\xi_1 + \xi_2$ is
ample (Nakai: positive square, positive on every effective class), and
general members of $|3(\xi_1+\xi_2)|$ are irreducible curves dominating both
factors. The point of (a)--(b) is that every such curve is \emph{numerically}
a fiber combination and acts as zero on $H^1$: as a correspondence it is
invisible. The collapse is a collapse of the correspondence calculus, not of
the surface's geometry --- and it is the calculus that the proposal required.
\end{remark}

\section{The diagonal surfaces}
\label{sec:diagonal}

On the diagonal surfaces $E_p\times E_p$, the endomorphism ring can be computed unconditionally.

\begin{theorem}[diagonal residue]
\label{thm:diagonal}
Fix a prime $p$.
\begin{enumerate}
\item $\End(E_p) = \Z$, unconditionally. (The lattice computation
below shows the only alternative would be $(\log p)^2 \in 4\pi^2\Q$ ---
which would make $\tau_p$ imaginary quadratic and $E_p$ a CM curve --- and
that coincidence is impossible: Gelfond--Schneider,
\S\ref{sec:transcendence}, Proposition. No $E_p$ has complex
multiplication.)
\item By (1), $\NS(E_p \times E_p) = \Z\xi_1 \oplus \Z\xi_2 \oplus \Z\cdot[\Delta]$
has rank $3$ and $\Corr(E_p, E_p) \cong \End(E_p) = \Z$, generated by the
diagonal. The available correspondences are (combinations of fiber classes
and) the graphs $\Gamma_m$ of multiplication by $m \in \Z$, with
$\Gamma_1 = \Delta$, and their complete intersection profile is:
\[
\Gamma_m\cdot\xi_1 = 1 , \quad
\Gamma_m\cdot\xi_2 = m^2 , \quad
(\Gamma_m\cdot\Gamma_n) = (m-n)^2 \ \ (m \neq n) , \quad
(\Gamma_m\cdot\Delta) = (m-1)^2
\]
(the last for all $m$: at $m = 1$ this is $\Delta^2 = 0$, by adjunction with
$K = 0$ and $e(E_p) = 0$).
\item Every number in this profile is a universal integer,
independent of $p$. More precisely: under the real-affine identification
$E_p \cong \R^2/\Z^2$ sending $(1, \tau_p)$ to the standard basis, the
subtori $\Gamma_m, \xi_1, \xi_2$ correspond to the \emph{same} subtori for
every $p$, and intersection numbers of divisor classes are topological (cup
products of integral classes). The map
\[
p \mapsto (\text{complete intersection profile of } E_p \times E_p)
\]
is constant. The prime enters the surface only through its complex structure
$\tau_p$, and the complex structure enters the correspondence calculus only
through $\Hom$/$\End$ --- which is the same ring $\Z$ for every $p$
(part~1).
\end{enumerate}
\end{theorem}

\begin{proof}
(1) $\lambda\Lambda_p \subseteq \Lambda_p$ with $\lambda = a + i b y_p$
(from $\lambda\cdot 1 \in \Lambda_p$) requires
$\lambda\cdot\tau_p = -b y_p^2 + i a y_p \in \Lambda_p$, i.e.\
$-b y_p^2 \in \Z$ (and the $\tau_p$-coefficient $a \in \Z$, automatic). If
$b \neq 0$ then $y_p^2 \in \Q$, i.e.\ $(\log p)^2 \in 4\pi^2\Q$ ---
impossible by the \S\ref{sec:transcendence} Proposition
(Gelfond--Schneider; a $b \neq 0$ endomorphism would make
$\tau_p^2 = -y_p^2 \in \Q$, i.e.\ $\tau_p$ quadratic imaginary and $\End$ a
CM order, and no $E_p$ admits this). So $b = 0$ and $\lambda = a \in \Z$:
$\End(E_p) = \Z$ with no exceptional case.

(2) Theorem~\ref{thm:NS} with $E = E' = E_p$ gives
$\NS = \Z\xi_1 \oplus \Z\xi_2 \oplus \End(E_p)$, and the graph of
$\mathrm{id}$ is $\Delta$. Intersection numbers: $\Gamma_m$ meets a vertical
fiber $\{x\}\times E_p$ in the single transverse point $(x, mx)$, so
$\Gamma_m\cdot\xi_1 = 1$. It meets a horizontal fiber $E_p\times\{c\}$ in
the fiber $[m]^{-1}(c)$, which has $m^2$ points, all transverse ($[m]$ is
unramified when $m\ne0$), so $\Gamma_m\cdot\xi_2 = m^2$ for $m\ne0$. For $m=0$, the graph is a horizontal fiber and its intersection with a distinct horizontal fiber is zero, giving the same formula. Here
$\#\ker([k]) = \#\bigl((\tfrac{1}{k}\Lambda_p)/\Lambda_p\bigr) = k^2$
(representatives $(i + j\cdot\tau_p)/k$, $0 \leq i, j < k$; check~D of the
numerical suite tabulates them for $k \leq 7$). For $m \neq n$,
$\Gamma_m \cap \Gamma_n = \{(x, mx) : (m-n)x = 0\} = \ker([m-n])$, with
$(m-n)^2$ points, each transverse (tangent directions $(1, m)$ and $(1, n)$
are independent), so $(\Gamma_m\cdot\Gamma_n) = (m-n)^2$; $n = 1$ gives
$(\Gamma_m\cdot\Delta) = (m-1)^2$ --- which is also the Lefschetz number
$1 - 2m + m^2$ of $[m]$ (trace on $H^0 - H^1 + H^2$), the integer Lefschetz
data of the family.

(3) The subtori $\Gamma_m$ are the images of $x \mapsto (x, mx)$, defined by
the \emph{integral} linear data $m$ alone; under the identification by the
basis $(1, \tau_p)$ they are the same integral subtori of $\R^4/\Z^4$ for
every $p$. Cup products of integral classes on a fixed topological manifold
are $p$-independent.
\end{proof}


\section{The precise connectedness obstruction}\label{sec:seed}
\begin{definition}
The cross-fiber graph has one vertex for each prime, and an undirected edge
$pq$ for $p\ne q$ exactly when $\Corr(E_p,E_q)\ne0$.
A connected classical cross-fiber system on a prime set $S$ means that this
graph, restricted to $S$, is connected. In particular, a proposed assembly
that requires a nonzero correspondence between every two distinct fibers
is a connected system in this sense.
\end{definition}

\begin{theorem}[matching obstruction]\label{thm:nogo}
Every connected component of the cross-fiber graph contains at most two
vertices. Hence a connected classical cross-fiber system on three or more
distinct prime fibers is impossible. In particular, an assignment requiring
nonzero classical correspondences for every pair of distinct primes cannot
exist.
\end{theorem}
\begin{proof}
Theorem~\ref{thm:NS} identifies the nonzero edge relation with
$\Hom(E_p,E_q)\ne0$. Theorem~\ref{thm:twopartner} says that each vertex has
degree at most one. A graph of maximum degree one has only singleton and
two-vertex connected components. This proves every assertion.
\end{proof}

\begin{remark}[what connectedness adds]
The hypothesis is about actual cross-fiber support. It is not implied by
merely asking for elements indexed by primes satisfying
$\Psi_m\circ\Psi_n=\Psi_{mn}$. In any diagonal ring
$\Corr(E_p,E_p)\cong\Z$, the classes of multiplication maps satisfy
$[m]\circ[n]=[mn]$ for positive integers $m,n$. These monoids exist regardless
of the vanishing of any cross-prime group. They do not by themselves realize
the Weil explicit formula, but they show why a multiplicative indexing
condition alone is insufficient for a connectedness obstruction.
\end{remark}

\section{Transcendence status}\label{sec:transcendence}
The diagonal and off-diagonal questions have different answers.
\begin{proposition}[no complex multiplication]
For every prime $p$, $(\log p)^2\notin4\pi^2\Q$. In particular
$\End(E_p)=\Z$ unconditionally.
\end{proposition}
\begin{proof}
If $(\log p)^2=4\pi^2r$, then $r\in\Q_{>0}$ and
$p=e^{2\pi\sqrt r}=(-1)^{-2i\sqrt r}$ using the branch $\log(-1)=i\pi$.
The base $-1$ is algebraic and different from $0,1$, and
$-2i\sqrt r$ is algebraic irrational. Gelfond--Schneider makes every value
of this power transcendental, contradicting that $p$ is an integer.
The lattice computation of Theorem~\ref{thm:diagonal} then gives the
endomorphism statement.
\end{proof}

For distinct primes, whether $\log p\log q\in4\pi^2\Q$ is possible remains
open. Its exclusion is the prime instance of the conjecture labeled
(C4E faible), attributed to Bertrand and printed by Diaz \cite[p.~231]{Dia97}:
for positive real algebraic $a,b\ne1$, $\pi^2$ and $\log a\log b$ are
$\Q$-linearly independent. The factor $4$ does not change the rational span.

For completeness the four-exponentials reduction is immediate. If
$\log p\log q=4\pi^2u/v$, take $x_1=2\pi i$, $x_2=\log p$,
$y_1=1$, $y_2=\log q/(2\pi i)$. Both pairs are linearly independent over
$\Q$, while the four exponentials are $1,q,p,e^{-2\pi iu/v}$, all algebraic.
This contradicts the four-exponentials conjecture \cite{FourExp}.
No proven six-exponentials theorem supplies the missing two-by-two
conclusion merely by renaming the input. None of the unconditional
connectedness or diagonal results assumes this conjecture.

\section{Consequences and limits for explicit-formula proposals}\label{sec:scope}
If an assembly requires nonzero classical divisorial correspondences joining
three distinct prime fibers, Theorem~\ref{thm:nogo} rules it out. On an
individual off-matching product the primitive divisor subgroup is zero and
no divisor acts nontrivially on $H^1$. These are exact restrictions on the
specified classical category.

On diagonal products, Theorem~\ref{thm:diagonal} gives a constant unmarked
intersection profile. For a finite real-coefficient test class,
\[
 D_{p,f}=\sum_m c_m(p,f)\Gamma_m+a(p,f)\xi_1+b(p,f)\xi_2,
 \qquad
 (D_{p,f}\cdot\Delta)=\sum_m c_m(p,f)(m-1)^2+a(p,f)+b(p,f).
\]
This formula identifies the information provided by the intersection lattice.
It does not prohibit dependence on $p$ or $f$ in the assignment. Indeed
$(\Gamma_p\cdot\Delta)=(p-1)^2$ varies with $p$ inside the same universal
profile. To exclude a test-family construction one must specify additional
restrictions on its coefficients and naturality, then prove an obstruction
under those restrictions.

Similarly, a prime-independent phase factor is not by itself an obstruction
to an archimedean contribution, which is a global term. The present results
neither construct nor rule out pole or archimedean terms in a more general
calculus. They concern classical divisor correspondences on the displayed
complex curves. A changed category, an enriched host, a disconnected
endomorphism construction, or additional analytic coefficient data requires
its own analysis. The characteristic-one constructions of Connes and Consani
and dynamical approaches are not excluded by this theorem.

\section*{Data, code and revision record}
The accompanying file \texttt{tate-products-v2-}\allowbreak\texttt{source-and-checks.zip}
contains this source, a portable exact-arithmetic check of the synthetic
isogeny index and graph intersection determinants, its output, build
instructions and the version-1-to-version-2 change record. The main results
are proved above; those finite checks are reproducibility diagnostics, not
proofs of the open off-diagonal transcendence conjecture. This revision omits
the optional decimal/PSLQ evidence because it is not needed by any theorem.
Historical version-1 artifacts remain at the publication's original record
and repository tag \texttt{paper-2026-08-28}. Cite the concept DOI
\href{https://doi.org/10.5281/zenodo.22171136}{10.5281/zenodo.22171136};
this revised manuscript is version 2.

\begin{thebibliography}{99}
\bibitem{CC26} A. Connes and C. Consani, \emph{On the Absolute Geometry
of Spec $\Z$ and the Fargues--Fontaine curve}, arXiv:2606.06604, 2026.

\bibitem{Wi04} J. Winkelmann, \emph{On elliptic curves in
$\mathrm{SL}_2(\C)/\Gamma$, Schanuel's conjecture and geodesic lengths},
Nagoya Math. J. \textbf{176} (2004), 159--180;
\href{https://doi.org/10.1017/S0027763000009016}{doi:10.1017/S0027763000009016};
arXiv:math/0204195v3. We use Lemma 7 and the proof of Theorem 2.

\bibitem{BL} C. Birkenhake and H. Lange, \emph{Complex Abelian Varieties},
2nd ed., Grundlehren 302, Springer, 2004, Chapters 2 and 5.

\bibitem{RS14} J. Rosen and A. Shnidman, \emph{N\'eron--Severi groups of
product abelian surfaces}, arXiv:1402.2233, 2014, Proposition 2.3.

\bibitem{Dia97} G. Diaz, \emph{La conjecture des quatre exponentielles
et les conjectures de D. Bertrand sur la fonction modulaire}, J. Th\'eor.
Nombres Bordeaux \textbf{9} (1997), 229--245, especially p.~231.

\bibitem{FourExp} S. Lang, \emph{Introduction to Transcendental Numbers},
Addison-Wesley, 1966; M. Waldschmidt, \emph{Diophantine Approximation on
Linear Algebraic Groups}, Grundlehren 326, Springer, 2000.
\end{thebibliography}
\end{document}
